\scanpage{414}
\mspage{\unsure{318|328}}
% archive sheet number: 413

% 1
fait que $J'$ est \unsure{grd|grand} devant \deleted{\unsure{Card}} $\operatorname{Card}\alpha=\operatorname{Card} I_\alpha$.

% 2–3
On doit de plus construire une flèche $u_\alpha:\varphi(i_\alpha)\longrightarrow\unsure{j|J}_\alpha$.

% 4–7
\deleted{\unsure{C}} C'est possible, car \deleted{il \unsure{exi}} $J'$ étant cofinale dans $J$, il existe $j'$ dans $J'$ et $\varphi(i_\alpha)\longrightarrow j'$, et $J'$ étant filtrante il existe $\overline{j_\alpha}$ et des flèches
% 8–9
\diagram{\begin{tikzcd}[ampersand replacement=\&]
j_\alpha\arrow[r] \& \overline{j_\alpha}\,;\\[-1ex]
j'\arrow[ur] \&
\end{tikzcd}}
% 9–12
donc quitte à remplacer $j_\alpha$ par $\overline{j}_\alpha$, les $\tau'_{\alpha,\beta}$ par des $\overline{\tau'_{\alpha,\beta}}$, on trouve bien une flèche \deleted{\unsure{$u_\alpha$}} $u_\alpha:\varphi(i_\alpha)\longrightarrow j_\alpha$. Il faut ensuite \deleted{vérifier les \unsure{cont}} s'occuper des compatibilités $3_\alpha(\beta)$, \deleted{Mais} donc des \deleted{\unsure{la}} doubles-flèches
% 13–15
\diagram{\begin{tikzcd}[ampersand replacement=\&,column sep=huge]
\varphi(i_\beta)\arrow[r,<-,shift left=2pt,"{\overline{\tau'}_{\alpha,\beta}u_\beta}","{\unsure{\,|\rangle}}"{at end,inner sep=0pt}]\arrow[r,<-,shift right=2pt,"{u_{\deleted{\unsure{\alpha}}\alpha}\varphi(\tau_{\alpha,\beta})}"'] \& j_\alpha\;.
\end{tikzcd}}

% 16–18
Comme \unsure{l'ens|l'ens.} de ces doubles flèches est de cardinal $\operatorname{Card}\alpha<\unsure{\pi|\varpi}$, et que $J^{\unsure{*|\prime}}$ est grd devant ledit cardinal, on
